\section{模型的评价与推广}

\subsection{模型的优点}

本文模型实现了任务瞬时资源占用与小时能源结算的分离。任务只要与某小时存在正重叠，其 GPU 和瞬时 IT 功率便被计入容量审计，而实际重叠时长被用于计算 GPU-hour 和电量。该处理能够同时适应分钟级任务和小时级电力数据，避免整小时计费造成的能源高估，也避免小时平均负荷掩盖瞬时容量冲突。

算力、设施与能源边界被置于统一可审计框架中。AI IT 负荷由任务类型功率映射和真实调度轨迹重建，非 AI IT 负荷保持不可迁移，PUE 作用于总 IT 负荷；新能源则被拆分为直接消纳、储能充电、外送和弃电。由此，成本、碳排放和新能源利用率均可依据任务与能源轨迹重新计算，而不依赖附件中的基准结果列。

模型还明确处理了时域末端和储能初始状态。第 $2400$ 至第 $2405$ 小时只承接此前到达的弹性任务，第 $2406$ 小时只用于终端结算；储能状态从第 $0$ 小时运行前的绝对初始 SOC 开始递推，并满足终端 SOC 约束。无储能比较状态与含储能构造保持逻辑独立，避免前者的可行性状态被错误传播至后者。

指标分母门禁和完整残差审计进一步提高了结果的可解释性。网络时延、新能源利用率和服务质量仅在相应分母有效时计算，不可用维度不会通过任意小常数强行进入评分。任务唯一执行、容量、功率、时延、截止时间、能源守恒、SOC 和互斥关系均在完整时域内接受复核。

跨情景结果揭示了不同外部因素的影响方向。售电价格系数为 $0$、$0.5$ 和 $1$ 时，运行成本分别为 $12327.7690$ 元、$-209670798.0048$ 元和 $-419353923.7785$ 元，而碳排放、网络时延、服务质量代理、新能源利用率和峰值净购电保持不变，表明该组情景的成本变化主要由售电收入核算驱动。峰谷价差系数为 $0.75$、$1$ 和 $1.25$ 时，成本分别为 $-419352397.8360$ 元、$-419353923.7785$ 元和 $-419355449.7210$ 元，其余报告指标保持不变。

新能源情景呈现出更显著的能源侧变化。新能源水平系数为 $0.8$ 时，成本、碳排放和峰值净购电分别为 $-302537839.3757$ 元、$30759.1226$ tCO$_2$ 和 $597.8863$ MW；系数为 $1.2$ 时，相应结果为 $-457183971.0773$ 元、$0$ tCO$_2$ 和 $0$ MW。固定种子 $2026$ 的新能源波动情景得到 $-402529632.4199$ 元的成本、$7923.8416$ tCO$_2$ 的碳排放和 $935.7721$ MW 的峰值净购电，说明新能源供给水平及波动会同时影响购售电收益、碳排放和电网峰值。

\subsection{模型的不足}

模型的主要不足在于数学模型与实际实现之间仍存在差距。需求预测采用固定 Huber 超参数，没有完成全部季节性候选及超参数组合的验证选模；任务调度没有执行契约规定的有限冲突回溯，问题二没有执行单任务局部改进，问题三没有执行储能动作候选搜索，问题四也没有执行基准与情景共享的大规模交替邻域搜索。因此，所得结果只能说明已构造候选通过全时域硬约束审计，无法排除未检查邻域中存在成本、碳排放或峰值表现更优的可行方案。

储能策略的工程刻画和动作机制均较为简化。规则策略仅利用新能源富余充电，累计充电量为 $2076.3155$ MWh，累计放电量为 $0$ MWh，电网充电量亦为 $0$；初始 SOC 合计为 $1594.0000$ MWh，终端 SOC 合计为 $3540.0000$ MWh。由于没有放电，储能不能减少购电量、峰值净购电或绝对爬坡量，也没有形成电价套利循环。

问题三的成本增加可由能源流分解得到直接解释。无储能与含储能状态的累计购电量均为 $31.5979$ MWh，而累计售电量分别为 $1330267.7686$ MWh 和 $1328783.5422$ MWh；相应运行成本分别为 $-419820327.7393$ 元和 $-419516103.8102$ 元。当前规则把部分原可外送新能源转化为终端库存，减少了售电收入，因而使成本增加 $304223.9291$ 元；购电量未改变，故碳排放保持为 $8.8506$ tCO$_2$。该结果只能说明现行只充不放规则未实现经济性与削峰改善，不能用于否定经过优化的储能调度。

成本指标采用“购电支出减售电收入”的货币量纲。问题二累计购电为 $40.4664$ MWh，累计售电为 $1329591.2956$ MWh，运行成本为 $-419658147.7076$ 元。负值表示在附件给定售电价格和区域最大外送边界下，累计售电收入超过购电支出，并不表示发电成本、储能成本或系统物理能耗为负。由于当前模型未纳入新能源建设、运维及储能退化等成本，该指标更准确地反映既定核算边界内的净购售电现金流。

通信模型只采用确定的单向网络时延，没有考虑带宽拥塞、迁移数据量、传输能耗和通信费用。服务质量采用 $1/(1+\text{平均等待小时})$ 的有限代理，而非按最大允许等待时间归一化的完整指标。对于网络资源紧张或任务迁移代价显著的实际系统，这些省略可能改变任务迁移方向和多目标权衡关系。

情景分析仍受固定任务方案限制。十五个情景均保持 $209$ 个迁移任务，网络时延为 $5.1594$ ms，服务质量代理为 $1.0$，说明情景变化主要通过能源流和储能轨迹复算体现，尚未反映任务迁移对价格、碳强度和新能源变化的响应。碳削减水平为 $0.1$、$0.2$ 和 $0.3$ 时，目标均在声明启发式内未达到；由于没有执行碳定向搜索，这一状态不构成数学不可行性证明。新能源波动分析也只采用固定种子 $2026$ 的单个情景，后续仍需在保留可复现种子的条件下开展重复仿真。

所有实际结果均被标记为 \texttt{non\_Pareto=true}，且没有全局或局部最优性证书。现有方案仅是确定性或固定种子启发式生成的 scenario-feasible 候选，不能被解释为 Pareto 前沿、局部最优解或连续与离散决策空间中的最优方案。

\subsection{模型的推广}

在算力调度方面，候选域方法可推广至多数据中心、边缘节点与云中心协同的任务分配问题。通过加入节点可用性、数据驻留、可信执行环境或任务亲和性约束，可处理具有隐私和合规要求的跨区域计算。对于允许抢占或分布式训练的任务，可将单一开工变量扩展为逐时执行份额，并增加任务连续性、检查点和同步通信约束。

在能源调度方面，区域独立平衡可扩展为含输电线路的网络模型。通过引入节点相角、线路容量或直流潮流约束，可研究任务迁移与电力跨区传输之间的替代关系。储能模型还可加入循环退化、温度、自放电、寿命成本和备用容量，使短期购售电收益、削峰需求与长期设备损耗得到统一权衡。

在不确定性处理方面，新能源出力、任务到达和执行时长可被构造为情景树或分布鲁棒集合。实际运行时，可在每个小时更新任务队列、负荷预测和新能源预测，并保留已开工任务的连续占用状态，从而形成可在线实施的模型预测控制框架\cite{rawlings2017model}。对于新能源波动，应采用多个预先固定的随机种子重复复算，并分别报告能源、碳排放与峰值指标的稳定性。

在求解方法方面，可先以本文确定性贪心策略生成可行初始方案，再采用混合整数规划、大邻域搜索或分解协调方法改进任务与储能决策。若计算预算有限，可按时延、截止裕度和最低资源下界筛除任务候选；若需要多目标决策，则可采用 $\varepsilon$-约束法生成经过验证的非支配候选集\cite{miettinen1999nonlinear}。无论采用何种扩展，任务轨迹、能源轨迹与评价指标均应在完整时域上重新计算，并继续保留结构化不可用状态、终止状态和约束残差审计。
